\section{Methodology}
\label{sec:method}

Figure~\ref{fig:method} outlines two stages: learning a support-conditioned semantic field from part annotations and using the frozen field to represent objects for manipulation.

\subsection{Problem Formulation}

At time $t$, a manipulation policy observes a scene point cloud $\mathcal{C}_t$, robot proprioception $q_t$, and segmented point clouds $\{\mathcal{P}_t^k\}_{k=1}^{K}$ for $K$ task-relevant objects, and predicts an action chunk $A_t$:
\begin{equation}
    A_t\sim\pi\!\left(\cdot\mid\mathcal{C}_t,q_t,
    \{R(\mathcal{P}_t^k)\}_{k=1}^{K}\right),
    \label{eq:policy_problem}
\end{equation}
where $R$ converts each object observation into a policy-facing representation. An object cloud is a partial measurement: changes in viewpoint, depth quality, segmentation, and sampling can alter its coordinates and attached features even when the object state is unchanged. We do not explicitly estimate these sensing factors. Instead, we seek an object representation that preserves 3D locations and functional-part information while reducing sensitivity to observation variation.

We realize $R$ with a support-conditioned semantic field. The observed object points form a support set $\mathcal{P}_{\mathrm{sup}}$, while an independently resampled set $\mathcal{Q}=\{x_j\}$ determines where policy-facing semantics are read:
\begin{equation}
    f_{\theta}(x\mid\mathcal{P}_{\mathrm{sup}})=(e_x,\ell_x),
    \qquad
    R(\mathcal{P}_{\mathrm{sup}})=\{(x,e_x):x\in\mathcal{Q}\},
    \label{eq:field}
\end{equation}
Here, $f_{\theta}$ denotes the complete field, with $\theta$ collecting all trainable parameters; $e_x\in\mathbb{R}^{d}$ is an L2-normalized embedding and $\ell_x\in\mathbb{R}^{C}$ contains training-time part logits. Queries need not be fixed across observations: each embedding is decoded from the support-conditioned field rather than attached to one support sample. We train a separate field per category on part-annotated PartNext models~\cite{wang2026partnext}, then freeze it for policy learning. The policy receives query coordinates and embeddings, but not part logits or labels.



\begin{figure*}[t]
    \centering
    \includegraphics[width=1\linewidth]{Images/method.png}
   % \vspace{-10pt}
    \caption{Overview of our object-centric continuous semantic field. Support points condition an object-specific tri-plane cache; 3D query locations read out part-aware embeddings and training-time part logits. The frozen field is then queried at resampled object locations to export semantic point clouds for downstream manipulation policies.}
    %\vspace{-15pt}
    \label{fig:method}
\end{figure*}

\subsection{Support-Conditioned Semantic Field}

Given an object observation $\mathcal{P}_{\mathrm{sup}}$, the field first constructs an object-conditioned cache and then decodes semantics at arbitrary queries. During field training, we sample $N_s=5{,}000$ support points uniformly with respect to mesh surface area and $N_q=2{,}048$ labeled surface queries using 75\% part-balanced and 25\% area-weighted sampling. At deployment, both sets instead come from the partial observed object cloud. In either case, we subtract the support centroid from both sets, divide their coordinates by the maximum support radius, and translate them so that the center of the support's $xy$ bounding box and its minimum $z$ coordinate lie at the origin.

We voxelize the normalized support at a grid size of $0.01$ and provide its coordinates, colors, and normals to a frozen Utonia encoder $E$~\cite{zhang2026utonia}. We use its final output restored to the input support resolution, without decoder-level feature upcasting. A lightweight learnable adapter $g_{\phi}$ maps each point feature into the field feature space:
\begin{equation}
    \widetilde{\mathcal{F}}_{\mathrm{sup}}
    =g_{\phi}\!\left(E(\mathcal{P}_{\mathrm{sup}})\right).
    \label{eq:support_encoding}
\end{equation}
The adapter consists of two linear layers with 256 channels, each followed by LayerNorm and GELU. It projects frozen encoder descriptors to the cache width and lets the field objectives adapt them to part semantics without updating Utonia. Its parameters $\phi$ belong to $\theta$. We then construct:
\begin{equation}
    \mathcal{T}_{\mathrm{obj}}
    =\Pi(\mathcal{P}_{\mathrm{sup}},\widetilde{\mathcal{F}}_{\mathrm{sup}})
    =\{T_{xy},T_{xz},T_{yz},F_{\mathrm{glob}}\},
    \label{eq:triplane}
\end{equation}
where $T_{xy}$, $T_{xz}$, and $T_{yz}$ are $64\!\times\!64$ feature maps with 256 channels, and $F_{\mathrm{glob}}\in\mathbb{R}^{256}$ is the mean adapted support feature. We normalize each axis to the support bounding box with 5\% padding, then bilinearly splat support features onto the three coordinate~planes.

For a query $x$, we use the same padded support bounds to project its coordinate onto each plane and bilinearly interpolate the corresponding features. The three 256-D samples are concatenated and fused by a linear--LayerNorm--GELU block into $F_{\mathrm{loc}}(x)\in\mathbb{R}^{256}$. A two-layer 256-D MLP receives the query coordinate, local feature, and replicated global feature, and predicts
\begin{equation}
    (e_x,\ell_x)=h_{\psi}\!\left(x,F_{\mathrm{loc}}(x),F_{\mathrm{glob}}\right).
    \label{eq:query_decoder}
\end{equation}
Here, $h_{\psi}$ is the query decoder within $f_{\theta}$, with $\psi\subset\theta$; the full field comprises support encoding, cache construction, and query decoding. The decoder produces a 128-D L2-normalized embedding and applies a linear classifier for part logits. Local features preserve spatial detail, while the global feature supplies instance context.


\subsection{Part-Aware Field Learning}
\label{sec:objectives}

The field must satisfy two complementary requirements: its embeddings should identify functional parts across object instances, and they should remain stable when the same object is observed differently. We therefore train it using paired augmented observations and three complementary losses.

\textbf{Paired observation construction.}
We transform each normalized support-query pair into two augmented observations:
\begin{equation}
    (\mathcal{P}_{\mathrm{sup}}^{(v)},\mathcal{Q}^{(v)})
    =\mathcal{T}_{\xi_v}(\mathcal{P}_{\mathrm{sup}},\mathcal{Q}),
    \qquad v\in\{1,2\},
    \label{eq:augmented_views}
\end{equation}
where $\mathcal{T}_{\xi_v}$ independently applies a uniform SO(3) rotation and Gaussian coordinate jitter with standard deviation $0.005$ to both sets. Each view is an augmented observation, not an additional camera. Shared surface samples and labels establish query correspondences despite transformed coordinates, allowing us to encourage semantic invariance. These augmentations model pose and local coordinate changes, but not occlusion, missing depth, or segmentation errors.

\textbf{Part anchoring.}
We apply part-label cross-entropy to the query logits in both views:
\begin{equation}
    \mathcal{L}_{\mathrm{ce}}
    =\frac{1}{2N_q}\sum_{v=1}^{2}\sum_{j=1}^{N_q}
    \mathrm{CE}(\ell_j^{(v)},y_j).
    \label{eq:part_loss}
\end{equation}
This loss supervises functional-part prediction at each query. To also encourage corresponding parts to occupy nearby locations in the embedding space used by the policy, we introduce a complementary alignment objective.

\textbf{Cross-instance semantic alignment.}
We therefore apply supervised contrastive learning independently in each view and average the results. For an embedding $e_i$, let $P(i)$ contain other queries in the minibatch with the same part label, including queries from different object instances. With temperature $\tau$,
\begin{equation}
\mathcal{L}_{\mathrm{con}}
=-\frac{1}{|\mathcal{I}|}\sum_{i\in\mathcal{I}}\frac{1}{|P(i)|}
\sum_{p\in P(i)}\log
\frac{\exp(e_i^\top e_p/\tau)}{\sum_{a\ne i}\exp(e_i^\top e_a/\tau)},
\label{eq:align_loss}
\end{equation}
where $\mathcal{I}$ contains anchors with at least one positive. We use $\tau=0.1$ and retain at most 2,048 valid queries per view. This loss pulls embeddings of the same functional part together and separates different parts, directly structuring the representation consumed by the policy.

\textbf{Paired-augmentation consistency.}
The contrastive loss aligns part-level clusters but does not explicitly require the same surface location to produce the same response under different observations. Using the known query correspondences, we minimize
\begin{equation}
    \mathcal{L}_{\mathrm{cr}}
    =\frac{1}{N_q}\sum_{j=1}^{N_q}
    \left(1-\langle e_j^{(1)},e_j^{(2)}\rangle\right).
    \label{eq:stability_loss}
\end{equation}
Because the embeddings are L2-normalized, this cosine loss directly penalizes semantic changes at corresponding locations. Unlike $\mathcal{L}_{\mathrm{con}}$, which establishes category-level part structure, $\mathcal{L}_{\mathrm{cr}}$ specifically reduces sensitivity to the modeled observation transform.

\textbf{Overall objective.}
We jointly optimize
\begin{equation}
    \mathcal{L}=\lambda_{\mathrm{ce}}\mathcal{L}_{\mathrm{ce}}
    +\lambda_{\mathrm{con}}\mathcal{L}_{\mathrm{con}}
    +\lambda_{\mathrm{cr}}\mathcal{L}_{\mathrm{cr}},
    \label{eq:total_loss}
\end{equation}
with $(\lambda_{\mathrm{ce}},\lambda_{\mathrm{con}},\lambda_{\mathrm{cr}})=(1.0,0.2,0.1)$. We train only the semantic adapter, local fusion block, and decoder for 4,000 epochs with minibatches of six object instances, using AdamW with learning rate $10^{-3}$ and weight decay $10^{-4}$. Utonia remains frozen, and training uses mixed precision. 



\subsection{Semantic Point Clouds for Manipulation}
\label{sec:method_policy}

After field training, all representation modules are frozen. At each observation frame, we independently sample 1,024 support points and $M=256$ query locations from the observed object cloud $\mathcal{P}_t^k$, using replacement only when fewer valid points are available. After joint normalization and field evaluation, we restore the query coordinates to the world frame and instantiate the representation in Eq.~\eqref{eq:policy_problem} as
\begin{equation}
    R(\mathcal{P}_t^k)=\mathcal{S}_{t}^{k}
    =\{(x_i,f_{\theta}^{\mathrm{emb}}(x_i\mid\mathcal{P}_{\mathrm{sup},t}^{k}))\}_{i=1}^{M}.
    \label{eq:semantic_cloud}
\end{equation}
The field is evaluated once per incoming observation. We use DP3~\cite{ze20243d} as the policy backbone. A 1,024-point scene cloud and each 131-channel semantic cloud (XYZ plus the 128-D embedding) are processed by separate max-pooled PointNet encoders, each producing a 128-D feature. These features are concatenated with proprioception and supplied to the unchanged diffusion objective and action representation.
